% Generated draft: requires mathematical and source review.
\documentclass[11pt]{article}
\usepackage[margin=25mm]{geometry}
\usepackage{fontspec}
\setmainfont{Latin Modern Roman}
\usepackage{amsmath,amssymb,mathtools,unicode-math}
\setmathfont{Latin Modern Math}
\usepackage{enumitem,needspace,xurl,hyperref}
\setlength{\parindent}{0pt}
\setlength{\parskip}{5pt}
\setlength{\emergencystretch}{3em}
\newcommand{\N}{\mathbb{N}}
\newcommand{\Z}{\mathbb{Z}}
\newcommand{\Q}{\mathbb{Q}}
\newcommand{\R}{\mathbb{R}}
\newcommand{\C}{\mathbb{C}}
\newcommand{\F}{\mathbb{F}}
\newcommand{\A}{\mathbb{A}}
\newcommand{\PP}{\mathbb{P}}
\newcommand{\E}{\mathbb{E}}
\newcommand{\eps}{\varepsilon}
\newcommand{\dd}{\,\mathrm{d}}
\begin{document}
\section*{UnsolvedMath ID: 30005397}
\textbf{Source problem code:} catalogue entry
\subsection{Definitions and mathematical statement}
The supplied title is \emph{Langlands--Rapoport Conjecture for Parahoric Level}. Write \(\mathrm{LR}\) and \(\mathrm{LR}\text{-}\tau\) only as names for the original Langlands--Rapoport conjecture and the variant mentioned in the extraction; their mathematical definitions are not supplied.

The extracted passage says that a proof of \(\mathrm{LR}\text{-}\tau\) ``should suffice for most applications'', although that conjecture ``is not as strong as the original LR conjecture''. This is a comparison and an expectation about applications, rather than a statement of either conjecture. The passage is attributed to Sug Woo Shin's contribution \emph{Cohomology of Shimura varieties}, reporting joint work with Mark Kisin and Yihang Zhu.

\textbf{Unresolved mathematical statement.} The supplied record does not specify a Shimura datum, a prime, hypotheses on the reductive group, the parahoric level, an integral model, or the objects and compatibility conditions involved in the proposed description. It also does not define the \(\tau\)-variant. Consequently, no precise quantified formulation of either conjecture, or of a parahoric-level specialization, can be recovered from this record. The title alone does not supply those hypotheses, and the extraction does not determine which conjecture is the intended target.

\textbf{Supplied status qualifications.} The registry labels this entry open, but also labels its statement unrecoverable. Its statement-verification note says that the conjecture discussed in the extraction has been proved and therefore is not an open problem. The supplied dated literature assessment distinguishes the two conjectures: it reports that \(\mathrm{LR}\text{-}\tau\) has been proved, with trace-formula consequences, while the original \(\mathrm{LR}\) conjecture remains wide open. These are supplied source and curation claims, not independently verified mathematical conclusions. Their distinction must be resolved before a precise open-problem statement can be assigned to this entry.
\subsection{Short English statement}
The passage compares the original Langlands--Rapoport conjecture with an unspecified \(\tau\)-variant; it does not state either conjecture. The supplied curation reports the variant proved and the original conjecture open. Missing definitions and hypotheses, together with conflicting registry qualifications, prevent recovery of a precise parahoric-level problem.
\subsection{Sources}
\begin{itemize}
\item Supplied local UnsolvedMath registry record, permanent ID \(30005397\), exact source problem code \texttt{catalogue entry}, in the Oberwolfach Reports Open Problems collection. Registry title: \emph{Langlands--Rapoport Conjecture for Parahoric Level}. The supplied proposed year is 2023; no proposer is specified.
\item \emph{Arithmetic of Shimura Varieties} (2023), Oberwolfach Reports, DOI 10.4171/owr/2023/5, pp.~261--326. Supplied report reference: \url{https://doi.org/10.4171/owr/2023/5}. The extraction attributes \emph{Cohomology of Shimura varieties} to Sug Woo Shin, joint work with Mark Kisin and Yihang Zhu. No precise page or theorem locator for either conjecture is supplied.
\item Supplied literature source: Kisin, Shin, and Zhu, \emph{Stable trace formula for Shimura varieties}, \url{https://arxiv.org/abs/2110.05381}.
\item Supplied literature source: Oberwolfach Report 5/2023, \url{https://doi.org/10.4171/owr/2023/5}.
\item Literature assessment checked at \texttt{2026-08-22}: the supplied curation reports that the LR-\(\tau\) conjecture has been proved and yields important trace-formula consequences, whereas the report describes the original Langlands--Rapoport conjecture as remaining wide open. This is supplied dated curation derived from the OWR material, not an independently verified assessment. The record requests independent verification of time-sensitive status and cited claims.
\end{itemize}
\end{document}
